\begin{workedexample}[Worked Example: 1~dB compression]
An amplifier has $20\ \text{dB}$ small-signal gain and an output 1~dB
compression point $\text{OP1dB} = +18\ \text{dBm}$. Driven at
$P_{in} = -5\ \text{dBm}$ the linear output would be $+15\ \text{dBm}$, below
OP1dB, so it is still roughly linear. Push toward $P_{in} = -1\ \text{dBm}$
(linear output $+19\ \text{dBm}$) and it compresses, clipping and generating
distortion.
\end{workedexample}

\begin{keytakeaways}
\begin{itemize}[leftmargin=1.4em,itemsep=2pt]
\item LNAs optimize noise figure; PAs optimize output power and efficiency; both need stability.
\item P1dB marks the onset of gain compression (the large-signal limit).
\item Higher linearity (IP3, P1dB) usually costs DC power and/or gain.
\end{itemize}
\end{keytakeaways}

\begin{exercises}
\item Gain $= 12\ \text{dB}$, $P_{in} = -30\ \text{dBm}$. What is the small-signal output?
\item Above P1dB, does the gain rise or fall?
\item Which amplifier goes first in a receive chain, and why?
\end{exercises}

\furtherreading{Gonzalez~\cite{gonzalez}; Razavi~\cite{razavi} (ch.~5).}
